\chapter{پاسخِ تمرین‌ها}%
\label{app:solutions}

پاسخ‌های پیشنهادی؛ اگر راهِ‌حلِ شما متفاوت است اما درست کار می‌کند، به
هیچ وجه نگران نباشید. در برنامه‌نویسی تقریباً همیشه بیش از یک راهِ درست
هست.

\section*{فصلِ ۱: برنامه‌نویسی یعنی چه؟}

\paragraph{۱.۱.} نمونه‌ها فراوان‌اند: ماشینِ لباس‌شویی (برنامهٔ شست‌وشو)،
تلویزیونِ هوشمند، آسانسور، عابربانک، ساعتِ هوشمند، جاروبرقیِ رباتیک، حتی
پلوپز. هر جا دستگاهی «تصمیم» می‌گیرد یا مرحله‌به‌مرحله کار می‌کند،
برنامه‌ای در کار است.

\paragraph{۱.۲.} نمونهٔ ممکن: ۱) مسواک را بردار. ۲) مسواک را خیس کن. % chktex 10
۳) به اندازهٔ یک نخود خمیردندان روی آن بگذار. ۴) هر دندان را از هر طرف % chktex 10
ده بار برس بزن. ۵) دهانت را با آب بشوی. ۶) مسواک را بشوی و سرِ جایش % chktex 10
بگذار. ابهام‌های رایجی که دوستتان می‌یابد: «هر طرف» یعنی دقیقاً کدام
سطح‌ها؟ «ده بار» با چه شدتی؟

\paragraph{۱.۳.} رایانه ۱۰ بار چاپ می‌کند؛ او دستور را اجرا می‌کند، نه
نیت را. اهمیتش این است: هر گاه برنامه خروجیِ «اشتباه» داد، اولین جایی که
باید بگردید دستورهای خودتان است.

\section*{فصلِ ۲: الگوریتم}

\paragraph{۲.۱.} کوچک‌ترینِ فعلی $\leftarrow$ عددِ اول؛ اگر عددِ دوم $<$
کوچک‌ترینِ فعلی آنگاه کوچک‌ترینِ فعلی $\leftarrow$ عددِ دوم؛ اعلام کن.
اجرای دستی: (۳ و ۸) جواب ۳؛ (۹ و ۹) جواب ۹، چون شرطِ «کوچک‌تر بودنِ اکید»
برقرار نمی‌شود و همان عددِ اول می‌ماند؛ که درست است.

\paragraph{۲.۲.} بله. با (۴ و ۴ و ۲): بزرگ‌ترینِ فعلی ۴ می‌شود؛ ۴ از ۴
بزرگ‌تر نیست، تغییری رخ نمی‌دهد؛ ۲ هم بزرگ‌تر نیست؛ جواب ۴ که درست است.
برابری آسیبی نمی‌زند چون شرطِ ما «بزرگ‌تر بودنِ اکید» است.

\paragraph{۲.۳.} باقی‌مانده $\leftarrow$ باقی‌ماندهٔ تقسیمِ عدد بر ۲؛
اگر باقی‌مانده $=$ ۰ آنگاه اعلام کن «زوج» وگرنه اعلام کن «فرد».

\paragraph{۲.۴.} گامِ «میانگین $\leftarrow$ مجموع $\div$ شمارنده» دچارِ
مشکل می‌شود: شمارنده صفر است و تقسیم بر صفر تعریف نشده. الگوریتمِ کامل
باید پیش از تقسیم بررسی کند شمارنده صفر نباشد (فصلِ ۶ همین را با
«نگهبان» حل کرد).

\paragraph{۲.۵.} تکرار: «به ازای هر رختِ روی بند\ldots»؛ انتخاب: «اگر چروک
بود اتو بزن»؛ توالی: جمع کردن، سپس اتو، سپس تا کردن، سپس در کمد گذاشتن.

\section*{فصلِ ۳: نخستین برنامه}

\paragraph{۳.۱.}
\begin{salamfa}
روال ریشه:
    سرچاپ "نامِ من: سارا"
    سرچاپ "شهر: شیراز"
پایان
\end{salamfa}

\paragraph{۳.۲.}
\begin{salamfa}
روال ریشه:
    سرچاپ "حاصل ضرب هفت در هشت:", 7 * 8
پایان
\end{salamfa}

\paragraph{۳.۳.}
\begin{salamfa}
روال ریشه:
    چاپ "یک، "
    چاپ "دو، "
    سرچاپ "سه!"
پایان
\end{salamfa}

\paragraph{۳.۴.} کامپایلر خطایی می‌دهد با این مضمون که روالِ ریشه
پایان نیافته و انتهای فایل زودتر از انتظار رسیده است؛ معمولاً به آخرین
خطِ فایل اشاره می‌کند. درس: هر \code{روال} (و بعدها هر \code{اگر} و
حلقه) یک \code{پایان} می‌خواهد.

\paragraph{۳.۵.} با نقل‌قول، عینِ متنِ \code{2 + 3} چاپ می‌شود؛
بی‌نقل‌قول، عبارت محاسبه می‌شود و \code{5} چاپ می‌شود.

\section*{فصلِ ۴: متغیرها و گونه‌های داده}

\paragraph{۴.۱.}
\begin{salamfa}
روال ریشه:
    نام کتاب : رشته = "شازده کوچولو"
    تعداد صفحه : صحیح = 96
    امتیاز : اعشار = 4.8
    سرچاپ نام کتاب, تعداد صفحه, امتیاز
پایان
\end{salamfa}

\paragraph{۴.۲.} خروجی \code{550} است: ۵۰۰ به‌علاوهٔ ۲۰۰ می‌شود ۷۰۰ و
منهای ۱۵۰ می‌شود ۵۵۰.
\begin{salamfa}
روال ریشه:
    ناپایا موجودی : صحیح = 500
    موجودی = موجودی + 200
    موجودی = موجودی - 150
    سرچاپ موجودی
پایان
\end{salamfa}

\paragraph{۴.۳.} چاپ می‌کند: \code{9 9}. با \code{آ = ب} مقدارِ ۴ برای
همیشه از دست می‌رود و هر دو جعبه ۹ می‌شوند؛ خطِ \code{ب = آ} هم دیگر
کاری نمی‌کند. جابه‌جایی شکست خورده است.

\paragraph{۴.۴.}
\begin{salamfa}
روال ریشه:
    ناپایا آ : صحیح = 4
    ناپایا ب : صحیح = 9
    موقت : صحیح = آ
    آ = ب
    ب = موقت
    سرچاپ آ, ب
پایان
\end{salamfa}
خروجی: \code{9 4}. همین ترفندِ «جعبهٔ موقت» در مرتب‌سازیِ فصلِ ۱۰ نقشِ
اصلی را بازی کرد.

\paragraph{۴.۵.} تعدادِ خواهر و برادر: صحیح (نصفه آدم نداریم!). وزن:
اعشار (۶۳٫۵ کیلوگرم معنا دارد). چراغ روشن است؟: منطقی (فقط دو حالت دارد).

\section*{فصلِ ۵: عملگرها}

\paragraph{۵.۱.} \code{10 - 4 * 2} می‌شود \code{2} (اول ضرب)؛
\code{(10 - 4) * 2} می‌شود \code{12}؛ \code{9 / 2} می‌شود \code{4}
(تقسیمِ صحیح)؛ \code{9.0 / 2.0} می‌شود \code{4.5}؛ \code{9 \% 2} می‌شود
\code{1}.

\paragraph{۵.۲.}
\begin{salamfa}
روال ریشه:
    دما : صحیح = 23
    مطبوع : منطقی = (دما >= 18) && (دما <= 28)
    سرچاپ "هوای مطبوع؟", مطبوع
پایان
\end{salamfa}

\paragraph{۵.۳.}
\begin{salamfa}
روال ریشه:
    عدد : صحیح = 2358
    سرچاپ "یکان:", عدد % 10
    سرچاپ "دهگان:", (عدد / 10) % 10
پایان
\end{salamfa}
\code{عدد / 10} در تقسیمِ صحیح می‌شود \code{235} و یکانِ آن، دهگانِ عددِ
اصلی است.

\paragraph{۵.۴.}
\begin{salamfa}
روال ریشه:
    سال : صحیح = 2000
    کبیسه : منطقی = (((سال % 4) == 0) && ((سال % 100) != 0)) ||
                    ((سال % 400) == 0)
    سرچاپ سال, کبیسه
پایان
\end{salamfa}
با 2024 حاصل \code{درست}، با 1900 \code{نادرست} (بر ۱۰۰ بخش‌پذیر است
ولی بر ۴۰۰ نه) و با 2000 \code{درست} است.

\paragraph{۵.۵.} \code{3 > 5} نادرست است؛ نقیضش درست؛ \code{2 == 2} هم
درست؛ و «درست \emph{و} درست» می‌شود \code{درست}.

\section*{فصلِ ۶: شرط‌ها}

\paragraph{۶.۱.}
\begin{salamfa}
روال ریشه:
    عدد : صحیح = -3
    اگر عدد > 0:
        سرچاپ "مثبت"
    وگرنه عدد < 0:
        سرچاپ "منفی"
    وگرنه:
        سرچاپ "صفر"
    پایان
پایان
\end{salamfa}

\paragraph{۶.۲.}
\begin{salamfa}
روال ریشه:
    الف : صحیح = 7
    ب : صحیح = 12
    ج : صحیح = 5
    ناپایا بزرگترین : صحیح = الف
    اگر ب > بزرگترین:
        بزرگترین = ب
    پایان
    اگر ج > بزرگترین:
        بزرگترین = ج
    پایان
    سرچاپ "بزرگ‌ترین:", بزرگترین
پایان
\end{salamfa}

\paragraph{۶.۳.} همان عبارتِ تمرینِ ۵.۴ درونِ یک \code{اگر}:
\begin{salamfa}
روال ریشه:
    سال : صحیح = 2024
    اگر (((سال % 4) == 0) && ((سال % 100) != 0)) ||
        ((سال % 400) == 0):
        سرچاپ سال, "کبیسه است"
    وگرنه:
        سرچاپ سال, "کبیسه نیست"
    پایان
پایان
\end{salamfa}

\paragraph{۶.۴.}
\begin{salamfa}
روال ریشه:
    الف : صحیح = 20
    ب : صحیح = 4
    عمل : رشته = "/"
    اگر عمل == "+":
        سرچاپ الف + ب
    وگرنه عمل == "-":
        سرچاپ الف - ب
    وگرنه عمل == "*":
        سرچاپ الف * ب
    وگرنه عمل == "/":
        اگر ب != 0:
            سرچاپ الف / ب
        وگرنه:
            سرچاپ "تقسیم بر صفر ممکن نیست"
        پایان
    وگرنه:
        سرچاپ "عمل ناشناخته"
    پایان
پایان
\end{salamfa}

\paragraph{۶.۵.} شرط‌ها وارونه چیده شده‌اند: با دمای ۵ چاپ می‌شود «هوا
سرد است» که اتفاقی درست به نظر می‌رسد، اما با دمای ۲۰ هم چاپ می‌شود
«هوا سرد است» (چون \code{20 > 0} همان اول درست است و به شرطِ معتدل
هرگز نمی‌رسیم) و با ۵- چاپ می‌شود «یخبندان!» که درست است. اصلاح: اول
سخت‌گیرترین شرط، یعنی \code{دما > 15} برای معتدل، بعد \code{دما > 0}
برای سرد و در پایان یخبندان.

\section*{فصلِ ۷: حلقه‌ها}

\paragraph{۷.۱.}
\begin{salamfa}
روال ریشه:
    ناپایا i : صحیح = 2
    تا i <= 20:
        سرچاپ i
        i = i + 2
    پایان
پایان
\end{salamfa}
با \code{تکرار}: \code{تکرار 1 تا 10 از i:} و درونش
\code{سرچاپ i * 2}.

\paragraph{۷.۲.}
\begin{salamfa}
روال ریشه:
    ناپایا i : صحیح = 10
    تا i >= 1:
        سرچاپ i
        i = i - 1
    پایان
    سرچاپ "پرتاب!"
پایان
\end{salamfa}

\paragraph{۷.۳.}
\begin{salamfa}
روال ریشه:
    ناپایا حاصل : صحیح = 1
    تکرار 1 تا 10 از i:
        حاصل = حاصل * i
    پایان
    سرچاپ "فاکتوریل ۱۰:", حاصل
پایان
\end{salamfa}
خروجی \code{3628800}. شروع از ۰ همه‌چیز را صفر می‌کرد، چون صفر ضرب در
هر چیزی صفر است؛ عنصرِ بی‌اثرِ ضرب ۱ است، همان‌طور که عنصرِ بی‌اثرِ جمع
۰ بود.

\paragraph{۷.۴.}
\begin{salamfa}
روال ریشه:
    ناپایا عدد : صحیح = 4725
    تا عدد > 0:
        سرچاپ عدد % 10
        عدد = عدد / 10
    پایان
پایان
\end{salamfa}
خروجی: ۵، بعد ۲، بعد ۷، بعد ۴.

\paragraph{۷.۵.}
\begin{salamfa}
روال ریشه:
    ناپایا سطر : صحیح = 4
    تا سطر >= 1:
        تکرار سطر:
            چاپ "*"
        پایان
        سرچاپ ""
        سطر = سطر - 1
    پایان
پایان
\end{salamfa}

\paragraph{۷.۶.} چهار بار. مقدارهای \code{i}: ۱۰، ۷، ۴، ۱ (هر چهار بار
شرط درست است و چاپ می‌شود)؛ سپس \code{i} به ۲- می‌رسد و حلقه می‌ایستد.

\section*{فصلِ ۸: روال‌ها}

\paragraph{۸.۱.}
\begin{salamfa}
روال کمینه(الف : صحیح, ب : صحیح) : صحیح:
    اگر الف < ب:
        برگشت الف
    پایان
    برگشت ب
پایان

روال ریشه:
    سرچاپ کمینه(3, 8)
    سرچاپ کمینه(40, 2)
    سرچاپ کمینه(6, 6)
پایان
\end{salamfa}

\paragraph{۸.۲.}
\begin{salamfa}
روال قدر مطلق(عدد : صحیح) : صحیح:
    اگر عدد < 0:
        برگشت -عدد
    پایان
    برگشت عدد
پایان
\end{salamfa}

\paragraph{۸.۳.}
\begin{salamfa}
روال زوج است(عدد : صحیح) : منطقی:
    برگشت (عدد % 2) == 0
پایان

روال ریشه:
    تکرار 1 تا 30 از i:
        اگر زوج است(i):
            سرچاپ i
        پایان
    پایان
پایان
\end{salamfa}

\paragraph{۸.۴.}
\begin{salamfa}
روال توان(پایه : صحیح, نما : صحیح) : صحیح:
    ناپایا نتیجه : صحیح = 1
    تکرار نما:
        نتیجه = نتیجه * پایه
    پایان
    برگشت نتیجه
پایان
\end{salamfa}

\paragraph{۸.۵.}
\begin{salamfa}
روال جداشمار مقسوم علیه ها(عدد : صحیح) : صحیح:
    ناپایا تعداد : صحیح = 0
    تکرار 1 تا عدد از ک:
        اگر (عدد % ک) == 0:
            تعداد = تعداد + 1
        پایان
    پایان
    برگشت تعداد
پایان

روال ریشه:
    تکرار 1 تا 30 از n:
        اگر جداشمار مقسوم علیه ها(n) == 2:
            سرچاپ n
        پایان
    پایان
پایان
\end{salamfa}
عددِ اول یعنی «فقط بر ۱ و خودش بخش‌پذیر»؛ یعنی دقیقاً دو مقسوم‌علیه.
عددِ ۱ فقط یک مقسوم‌علیه دارد و به همین دلیل اول به شمار نمی‌آید؛ این
تعریف آن حالتِ مرزی را خودبه‌خود درست از آب درمی‌آورد.

\section*{فصلِ ۹: آرایه‌ها}

\paragraph{۹.۱.}
\begin{salamfa}
روال ریشه:
    دماها : صحیح[7] = [21, 24, 19, 28, 31, 22, 25]
    ناپایا مجموع : صحیح = 0
    هر دما از دماها:
        مجموع = مجموع + دما
    پایان
    میانگین : اعشار = (مجموع * 1.0) / len(دماها)
    سرچاپ "میانگین:", میانگین
    ناپایا تعداد : صحیح = 0
    هر دما از دماها:
        اگر دما > میانگین:
            تعداد = تعداد + 1
        پایان
    پایان
    سرچاپ "روزهای گرم‌تر از میانگین:", تعداد
پایان
\end{salamfa}
دو پیمایش لازم است چون تا پیمایشِ اول تمام نشود، میانگین معلوم نیست و
بدونِ میانگین نمی‌توان مقایسه کرد.

\paragraph{۹.۲.} همان \code{بیشینه آرایه} است با وارونه کردنِ جهتِ
مقایسه: \code{اگر عدد < کوچکترین:}. (متنِ کاملش در پروژهٔ پایانیِ فصلِ
۱۱ آمده است.)

\paragraph{۹.۳.}
\begin{salamfa}
روال ریشه:
    ناپایا مربع ها : صحیح[10] = [0, 0, 0, 0, 0, 0, 0, 0, 0, 0]
    تکرار 10 از i:
        مربع ها[i] = i * i
    پایان
    هر مقدار از مربع ها:
        سرچاپ مقدار
    پایان
پایان
\end{salamfa}

\paragraph{۹.۴.}
\begin{salamfa}
روال ریشه:
    اعداد : صحیح[5] = [10, 20, 30, 40, 50]
    ناپایا i : صحیح = len(اعداد) - 1
    تا i >= 0:
        سرچاپ اعداد[i]
        i = i - 1
    پایان
پایان
\end{salamfa}
اولین اندیس \code{len - 1} است (نه \code{len}!) و آخرین اندیس صفر؛ پس
شرط \code{i >= 0} است، نه \code{i > 0}. هر دو سرِ حلقه جای خطای یکی کم
و زیاد بود.

\paragraph{۹.۵.}
\begin{salamfa}
روال جداشمار(اعداد : صحیح[:], هدف : صحیح) : صحیح:
    ناپایا تعداد : صحیح = 0
    هر عدد از اعداد:
        اگر عدد == هدف:
            تعداد = تعداد + 1
        پایان
    پایان
    برگشت تعداد
پایان

روال ریشه:
    داده : صحیح[5] = [3, 7, 3, 3, 1]
    سرچاپ جداشمار(داده, 3)
پایان
\end{salamfa}

\section*{فصلِ ۱۰: الگوریتم‌های کلاسیک}

\paragraph{۱۰.۱.} راهِ «پیمایش از آخر»:
\begin{salamfa}
روال جستجوی آخرین(اعداد : صحیح[:], هدف : صحیح) : صحیح:
    ناپایا i : صحیح = len(اعداد) - 1
    تا i >= 0:
        اگر اعداد[i] == هدف:
            برگشت i
        پایان
        i = i - 1
    پایان
    برگشت -1
پایان
\end{salamfa}
راهِ دوم: از اول پیمایش کنید، هر بار که هدف را دیدید اندیس را در متغیری
مانندِ \code{آخرین دیده شده} (با مقدارِ اولیهٔ \code{-1}) بنویسید و در
پایان همان را برگردانید؛ این راه \code{بازگشتِ} زودهنگام ندارد.

\paragraph{۱۰.۲.} آرایهٔ \code{[5, 1, 4, 2]}: پس از گشتِ اول
\code{[1, 4, 2, 5]} (سه مقایسه، سه جابه‌جایی: ۵ تا انتها حباب می‌زند)؛
پس از گشتِ دوم \code{[1, 2, 4, 5]}؛ گشتِ سوم چیزی را عوض نمی‌کند؛ مرتب
است.

\paragraph{۱۰.۳.} فقط جهتِ مقایسه: \code{اگر اعداد[i] < اعداد[i + 1]:}
(علامتِ \code{>} به \code{<} تبدیل می‌شود).

\paragraph{۱۰.۴.}
\begin{salamfa}
روال ک م م(الف : صحیح, ب : صحیح) : صحیح:
    برگشت (الف / ب م م(الف, ب)) * ب
پایان
\end{salamfa}
اول تقسیم و بعد ضرب می‌کنیم تا عددِ میانی بی‌جهت بزرگ نشود؛ تقسیم هم
دقیق است چون ب.م.م، مقسوم‌علیهِ \code{الف} است.

\paragraph{۱۰.۵.}
\begin{salamfa}
روال مجموع رقم ها(عدد : صحیح) : صحیح:
    ناپایا مانده : صحیح = عدد
    ناپایا مجموع : صحیح = 0
    تا مانده > 0:
        مجموع = مجموع + (مانده % 10)
        مانده = مانده / 10
    پایان
    برگشت مجموع
پایان
\end{salamfa}

\paragraph{۱۰.۶.}
\begin{salamfa}
روال جستجوی دودویی(اعداد : صحیح[:], هدف : صحیح) : صحیح:
    ناپایا پایین : صحیح = 0
    ناپایا بالا : صحیح = len(اعداد) - 1
    تا پایین <= بالا:
        وسط : صحیح = (پایین + بالا) / 2
        اگر اعداد[وسط] == هدف:
            برگشت وسط
        وگرنه اعداد[وسط] < هدف:
            پایین = وسط + 1
        وگرنه:
            بالا = وسط - 1
        پایان
    پایان
    برگشت -1
پایان
\end{salamfa}
دقت کنید که \code{وسط} در هر دورِ حلقه دوباره حساب می‌شود و «کوچک‌تر
بودنِ خانهٔ وسط از هدف» یعنی جواب اگر باشد در نیمهٔ بالایی است.

\section*{فصلِ ۱۱: پروژهٔ پایانی}

\paragraph{۱۱.۱.} یک خط در روالِ ریشه کافی است:
\code{سرچاپ "فاصله نمره‌ها:", بیشینه آرایه(نمرات) - کمینه آرایه(نمرات)}.
(اگر پس از مرتب‌سازی هستید، تفاضلِ خانهٔ آخر و اول هم همان است.)

\paragraph{۱۱.۲.} همان قالبِ \code{شمارش قبولی} با شرطِ
\code{نمره < 10} یعنی نمرهٔ تک‌رقمی؛ یا سخت‌گیرانه‌تر،
\code{(نمره >= 0) \&\& (نمره <= 9)}.

\paragraph{۱۱.۳.}
\begin{salamfa}
روال ریشه:
    تکرار 100 تا 999 از n:
        اگر (n == وارونه(n)) && اول است(n):
            سرچاپ n
        پایان
    پایان
پایان
\end{salamfa}
(روال‌های \code{وارونه} و \code{اول است} را از فصل‌های ۸ و ۱۰ بیاورید.)
خروجی با ۱۰۱ شروع می‌شود و اعضایی مانندِ ۱۳۱، ۱۵۱، ۱۸۱ و ۱۹۱ دارد.

\paragraph{۱۱.۴.}
\begin{salamfa}
روال مرتب سازی انتخابی(اعداد : صحیح[:]):
    ن : صحیح = len(اعداد)
    ناپایا i : صحیح = 0
    تا i < (ن - 1):
        ناپایا جای کمینه : صحیح = i
        ناپایا j : صحیح = i + 1
        تا j < ن:
            اگر اعداد[j] < اعداد[جای کمینه]:
                جای کمینه = j
            پایان
            j = j + 1
        پایان
        اگر جای کمینه != i:
            موقت : صحیح = اعداد[i]
            اعداد[i] = اعداد[جای کمینه]
            اعداد[جای کمینه] = موقت
        پایان
        i = i + 1
    پایان
پایان
\end{salamfa}
خروجی همان خروجیِ مرتب‌سازیِ حبابی است؛ تفاوت در راهِ رسیدن است:
انتخابی در هر دورِ بیرونی فقط \emph{یک} جابه‌جایی می‌کند (پیدا کن، بعد
جابه‌جا کن)، در حالی که حبابی ممکن است در یک گشت بارها جابه‌جا کند.
